% Figure: level sets q(x)=1 of three quadratic forms, showing how the
% eigenvalue signs set the geometry: positive definite (ellipse),
% indefinite (hyperbola), and positive semidefinite (parallel lines).
\begin{figure}[htbp]
\centering
\begin{tikzpicture}[scale=1.0]
% Panel 1: positive definite, ellipse.
\begin{scope}[shift={(0,0)}]
\draw[->,gray] (-2.2,0) -- (2.2,0) node[right,font=\scriptsize]{$x_1$};
\draw[->,gray] (0,-2.2) -- (0,2.2) node[above,font=\scriptsize]{$x_2$};
% q = 9 y1^2 + y2^2 with y1=(x1+x2)/sqrt2, y2=(x1-x2)/sqrt2; the level
% set q=1 is an ellipse tilted 45 deg, semi-axes 1/3 and 1 along the
% eigenvectors (1,1)/sqrt2 and (1,-1)/sqrt2.
\draw[engblue, very thick, rotate=45] (0,0) ellipse [x radius=0.3333, y radius=1];
\draw[engred, ->, thick] (0,0) -- (0.2357,0.2357);
\draw[engamber, ->, thick] (0,0) -- (0.7071,-0.7071);
\node[font=\scriptsize, engblue, anchor=north] at (0,-2.5) {positive definite};
\node[font=\scriptsize, anchor=north] at (0,-2.95)
  {$\lambda_1,\lambda_2 > 0$: ellipse};
\end{scope}
% Panel 2: indefinite, hyperbola.
\begin{scope}[shift={(6,0)}]
\draw[->,gray] (-2.2,0) -- (2.2,0) node[right,font=\scriptsize]{$x_1$};
\draw[->,gray] (0,-2.2) -- (0,2.2) node[above,font=\scriptsize]{$x_2$};
% q = x1^2 - x2^2 = 1: hyperbola opening along x1.
\draw[engblue, very thick, domain=0:1.317, samples=40, variable=\t]
  plot ({cosh(\t)},{sinh(\t)});
\draw[engblue, very thick, domain=0:1.317, samples=40, variable=\t]
  plot ({cosh(\t)},{-sinh(\t)});
\draw[engblue, very thick, domain=0:1.317, samples=40, variable=\t]
  plot ({-cosh(\t)},{sinh(\t)});
\draw[engblue, very thick, domain=0:1.317, samples=40, variable=\t]
  plot ({-cosh(\t)},{-sinh(\t)});
% asymptotes.
\draw[gray, dashed] (-2,-2) -- (2,2);
\draw[gray, dashed] (-2,2) -- (2,-2);
\node[font=\scriptsize, engblue, anchor=north] at (0,-2.5) {indefinite};
\node[font=\scriptsize, anchor=north] at (0,-2.95)
  {$\lambda_1 > 0 > \lambda_2$: hyperbola};
\end{scope}
% Panel 3: positive semidefinite, degenerate.
\begin{scope}[shift={(12,0)}]
\draw[->,gray] (-2.2,0) -- (2.2,0) node[right,font=\scriptsize]{$x_1$};
\draw[->,gray] (0,-2.2) -- (0,2.2) node[above,font=\scriptsize]{$x_2$};
% q = x1^2 = 1: two parallel lines x1 = +-1 (one zero eigenvalue).
\draw[engblue, very thick] (1,-2) -- (1,2);
\draw[engblue, very thick] (-1,-2) -- (-1,2);
\node[font=\scriptsize, engblue, anchor=north] at (0,-2.5) {semidefinite};
\node[font=\scriptsize, anchor=north] at (0,-2.95)
  {$\lambda_1 > 0 = \lambda_2$: parallel lines};
\end{scope}
\end{tikzpicture}
\caption[Level sets of quadratic forms by eigenvalue sign]{The level
set $q(\vect{x}) = \vect{x}\transpose\mat{A}\vect{x} = 1$ for three
symmetric matrices, one per panel, with the geometry decided entirely
by the signs of the eigenvalues. When both eigenvalues are positive
(left) the level set is an ellipse whose semi-axes are
$1/\sqrt{\lambda_i}$ along the eigenvectors (red and amber arrows show
the principal directions of the tilted example
$\protect\begin{psmallmatrix}5\&4\\4\&5\protect\end{psmallmatrix}$,
nine times stiffer along one diagonal). When the eigenvalues straddle
zero (centre) the form is indefinite and the level set is a hyperbola
with the dashed asymptotes running along the directions of zero form
value. When one eigenvalue vanishes (right) the form is positive
semidefinite and the ellipse degenerates into a pair of parallel
lines, the level set flat along the null direction. The same sign
reading classifies a Hessian at a critical point: bowl, saddle, and
trough.}
\label{fig:vol02:ch10:quadratic-conics}
\end{figure}
